Omdat de log-bestanden plain-text kunnen we alle standaard teksttools op deze bestanden loslaten.

Om direct te kunnen volgen wat er op een systeem gebeurt wordt veel het \texttt{tail} commando gebruikt. \texttt{tail} heeft een optie \texttt{-f} die follow betekent. Het blijft volgen welke nieuwe regels er aan een bestand worden toegevoegd. Een manier om bijvoorbeeld de \texttt{messages} log te blijven volgen is:
\begin{lstlisting}[language=bash]
$ sudo tail -f /var/log/messages
\end{lstlisting}

Omdat alle log-bestanden tekst bestanden zijn kan je ook bijvoorbeeld \texttt{grep} erop los laten en zoeken naar bepaalde patronen in de log-bestanden.

